\chapter{ADMET}
La valutazione delle proprietà \textbf{ADMET} (Assorbimento, Distribuzione, Metabolismo, Escrezione e Tossicità) è considerata un passaggio critico all'interno della pipeline di svilluppo farmacologico, al fine di ridurre "a monte" i tassi di fallimento nelle fasi cliniche. Infatti, la scorretta valutazione di queste proprietà rappresenta un incremento nel rischio di logoramento dei farmaci candidati.

\section{Il framework ADMET}
L'impiego del Machine Learning in questo campo rappresenta un'importante svolta metodologica rispetto ai test empirici tradizionali \cite{Yang2021}, poiché permette di processare grandi volumi di dati molecolari per prevederne il comportamento farmacocinetico con costi e tempi drasticamente ridotti, allargando il primo "collo di bottiglia" del processo di ricerca.

In questo nuovo contesto, il framework ADMET non viene più interpretato come una serie di test isolati, ma come un sistema integrato di vincoli biochimici che si influenzano a vicenda e che determinano la sicurezza di una molecola. 

\section{La tossicità molecolare}
TO-DO

\section{Evoluzione verso il Machine Learning per la predizione ADMET} %da li2026
Storicamente, i primi tentatici di razionalizzare le proprietà e l'attività biologica delle molecole si basava su metodi come \textit{quantitative structure-activity relationship} (QSAR) e \textit{structure-activity relationship} (SAR).
Tuttavia questi metodi presentavano un limite fondamentale: la necessità di una selezione manuale e di un'accurata rappresentazione delle \textit{feature} chimiche. %Approfondisci
Con l'avvento del ML classico, tuttavia, l'attenzione si è spostata verso l'uso di rappresentazioni come i \textit{chemical fingerprints}, che vengono presi come input da algoritmi come \textbf{Random Forest}, \textbf{Support Vector Machines}, \textbf{k-Nearest Neighbors} e \textbf{XGBoost}. 
Questi ultimi hanno permesso di apprendere classificazioni complesse, trattando la molecola come un vettore binario che indica presenza/assenza di frammenti chimici, e raggiungendo risultati in modo veloce e con processi parzialmente interpretabili.
Tuttavia, questi algoritmi ignorano la struttura tridimensionale e la connettività delle molecole, e le leggono invece come una sequenza di caratteristiche booleane, rendendo le interazioni tra frammenti vicini nello spazio ma lontani nella stringa di difficile predizione.

Per una predizione più accurata di queste interazioni più complesse, \textbf{Li et al. (2026)} propongono un modello chiamato \textbf{CMA}, ovvero "Cross-Aligned Multimodal Attention". 
L'idea centrale è proprio quella di rappresentare la molecola come in un \textbf{attention cross-modale}, ovvero tramite tre le matrici Query (Q), Key (K) e Value (V) del meccanismo di attentio dei Transformer. Ogni matriche viene mappata a partire da una rappresentazione diversa, e solo una di queste fa riferimento ai \textit{fingerprints}:
\begin{itemize}
    \item Query utilizza \textbf{GROVER}, un algoritmo che decodifica informazioni sulla struttura delle molecole in un grafo in cui ogni nodo rappresenta un atomo e i legami chimici sono astratti come archi. GROVER è un \textbf{GNN} (Graph Neural Networks) preaddestrato su 10 milioni di molecole e permette la la rappresentazione più ricca strutturalmente. Contruibuisce a circa il 36\% della contibuzione media finale.
    \item Value viene processata da \textbf{ResNet-50}, una CNN in grado di estrarre immagini molecolari bidimensionali e catturare così la forma, la simmetria e l'arrangiamento degli atomi dei soggetti candidati. La rete viene preaddestrata su 10 milioni di molecole. Contruibuisce a circa il 22\% della contibuzione media finale.
    \item Key, infine, fa uso di \textbf{fingerprint MACCS}, composta da 166 bit che codifica l'occorrenza di gruppi chimici specifici. Contruibuisce a circa il 42\% della contibuzione media finale.
\end{itemize}
La funzione di punteggio di Attention è definita come segue:
$$\text{Attention}(Q, K, V) = \text{softmax}\left(\frac{QK^T}{\sqrt{d_k}}\right)V$$
La modalità Query "interroga" la modalità Key per determinare quali elementi siano più rilevanti all'interno del contesto corrente. Il risultato è infine un media pesata dalla modalità Value. Tradotto nell'ambito molecolare, la rappresentazione a grafo può interrogare le \textit{fingerprint} per comprendere al meglio quali siano le caratteristiche strutturali locali più coerenti con la topologia generale, mentre il risultato finale viene modulato dalle immagini estratte da ResNet. 
%TO_DO interpretabilità e gradcam
\cite{li2026}

\section{Multimodalità e meccanismi di Attention}
%\section{Modelli predittivi/GNN per la previsione della sicurezza}